\newcommand{\stationpanel}[4]{\begin{tikzpicture}
\begin{axis}[width=#1\linewidth,height=4.65cm,xmin=#2,xmax=#3,ymin=.25,ymax=.80,
title={#4},title style={font=\sffamily\small\bfseries,color=ink},
axis lines*=left,ytick={.30,.38,.60,.76},scaled ticks=false,ymajorgrids,grid style={ink!13},
tick label style={font=\sffamily\scriptsize},xlabel={时间 (s)},xlabel style={font=\sffamily\footnotesize}]
\addplot[ink!75,densely dashed] coordinates {(#2,.38) (#3,.38)};
\addplot[teal,thick,no marks] table[col sep=comma,x=time_s,y=fds_m]{data/pinn-end-station.csv};
\addplot[rust,thick,dashed,no marks] table[col sep=comma,x=time_s,y=pinn_m]{data/pinn-end-station.csv};
\end{axis}\end{tikzpicture}}
\reportfigure{{\sffamily\small 烟层高度 (m)\hfill\textcolor{teal}{\rule[.4ex]{12pt}{1.1pt}} FDS\qquad\textcolor{rust}{-- --} PINN\quad\textcolor{gray}{-- --} 0.38 m 阈值\par}
\stationpanel{.57}{0}{120}{完整时序}\hfill\stationpanel{.40}{81}{85}{越限时段}}{图 E2.\quad 留出的 OP-6 工况，垂壁深度 0.24~m、烟炱产率 0.05。两图显示远端 $x=6.1$~m 同一组原始采样值。83.0~s 时 FDS 记录为 0.353~m，PINN 预测为 0.750~m。八个测点的全时段结果见图 E1。}
